% ============================================================
% Section 2 — Related Work (body-only; \input into root.tex)
% Facts: notes/fact_related.md (taxonomy plan ii, contrasts, per-competitor table iii)
% All \cite keys exist in latex/references.bib.
% ============================================================

\section{Related Work}
\label{sec:related}

\subsection{Multimodal Semantic Segmentation}
\label{sec:related:mmss}

Benchmarks with synchronized heterogeneous sensors have driven this field: DELIVER~\cite{cmnext2023} pairs RGB with depth, event, and LiDAR data under five adverse conditions and corruption cases, MUSES~\cite{muses2024} provides multi-sensor driving scenes under visual degradation, and MCubeS~\cite{mcubes2022} targets material segmentation. On the architecture side, CMX~\cite{cmx2023} introduced cross-modal feature rectification and fusion for RGB-X segmentation, TokenFusion~\cite{tokenfusion2022} substitutes uninformative tokens across modalities inside a ViT, GeminiFusion~\cite{geminifusion2024} performs pixel-wise intra- and inter-modal attention, and StitchFusion~\cite{stitchfusion2024} weaves features between frozen pretrained encoders through lightweight adapters. A second line pursues \emph{arbitrary-modal} robustness: MAGIC and MAGIC++~\cite{magic2024,magicpp2024} counter RGB-centric bias with hierarchical modality selection, AnySeg~\cite{anyseg2024} distills uni- and cross-modal knowledge to survive missing sensors, and OmniSegmentor~\cite{omnisegmentor2025} pretrains a single encoder on a five-modality ImageNeXt corpus, reaching 68.0 mIoU on the DELIVER \emph{val} split---a pretraining axis that is orthogonal to, and composable with, our adaptation-based approach. Common to all of these methods is that they adapt \emph{representations or architectures} to multiple sensors; none of them decides at inference time \emph{how much each sensor should be trusted}.

\subsection{Condition- and Reliability-Aware Fusion}
\label{sec:related:reliability}

A more recent family conditions the fusion process on scene state or sensor quality. CAFuser~\cite{cafuser2025} derives a condition token from the RGB stream in CLIP~\cite{clip2021} space and uses it to steer fusion, which requires verbo-visual contrastive supervision and condition annotations. DGFusion~\cite{dgfusion2026}, the current DELIVER \emph{test} state of the art (56.7 mIoU, four-modality CLDE setting), treats LiDAR both as an input and as a supervision target: ground-truth depth trained with a robust log-depth loss yields local depth tokens and a global condition token that condition cross-modal fusion. HyperDUM~\cite{hyperdum2025} estimates deterministic uncertainty from hyperdimensional prototype distances---prototypes built from labels plus a fine-tuned weighting layer---and reweights features before fusion (66.30$\rightarrow$67.59 mIoU on DELIVER \emph{val}). Evidential approaches learn uncertainty heads: TMC~\cite{tmc2021} aggregates Dirichlet opinions for multi-view classification and UTFNet~\cite{utfnet2023} applies uncertainty-guided feature weighting to RGB-thermal segmentation. Reliability has also been injected at the loss or adaptation level: functional-entropy regularization balances modality contributions during training only~\cite{zheng2025entropy}, READ~\cite{read2024} reweights a test-time-adaptation loss against modality reliability bias in classification, and ReliFusion~\cite{relifusion2025} scales the \emph{output} of cross-attention by a learned reliability score for 3D detection. Across this family, the reliability or condition signal invariably (i) depends on supervision beyond segmentation labels---ground-truth depth, condition text or annotations, label-built prototypes, or dedicated evidential losses---and/or (ii) enters at the feature, attention-output, or loss level rather than where per-class decisions are formed.

\subsection{Foundation Models and Parameter-Efficient Adaptation}
\label{sec:related:vfm}

Parameter-efficient fine-tuning makes frozen vision foundation models practical for dense prediction: LoRA~\cite{lora2022} adds low-rank updates to frozen projections, ViT-Adapter~\cite{vitadapter2023} equips plain ViTs with dense-prediction inductive biases, and SAMed~\cite{samed2023} established the LoRA-on-SAM precedent. Building on SAM and SAM~2~\cite{sam2023,sam2_2024}, MemorySAM~\cite{memorysam2025} treats modalities as video ``frames'' fused by SAM~2 memory attention---the architectural ancestor of the attention-level reliability biasing that we audit and ultimately reject in Sec.~\ref{sec:experiments}. MLE-SAM~\cite{mlesam} attaches modality-specific mixture-of-LoRA experts with a learned dispatch gate to a frozen SAM encoder, and MM~SAM-adapter~\cite{mmsamadapter2025} injects fused features into a SAM ViT-L via deformable cross-attention, but supports only two modalities per model.\footnote{MM~SAM-adapter reports DELIVER test results under its own two-modality protocol (e.g., 57.35 mIoU RGB-D), which is not comparable to the four-modality CLDE setting used in Table~\ref{tab:sota}.} In contrast to the SAM lineage, DINOv3~\cite{dinov3_2025} provides a self-supervised backbone with strong frozen dense features; to our knowledge, as of mid-2026, no prior multimodal segmentation method builds on DINOv3.\todo{re-run the DINOv3-multimodal-seg arXiv sweep at submission time; nearest prior work uses DINOv2 with a click-based NoC metric rather than mIoU} These adaptation methods answer \emph{how to represent} each modality on a frozen backbone, but not \emph{how much to trust} it at a given pixel and class.

\medskip
\noindent\textbf{Positioning.}
ReliaDINO combines both threads under a minimal supervision budget: a single frozen DINOv3 ViT-L multiplexed by per-modality LoRA adapters, with reliability applied \emph{only} where class decisions are formed---a calibrated competence gate at the fusion output and a per-class reliability-anchored router---rather than as an attention-output scale~\cite{relifusion2025}, a TTA loss weight~\cite{read2024}, or a feature reweighting stage~\cite{hyperdum2025}. Unlike DGFusion~\cite{dgfusion2026} and CAFuser~\cite{cafuser2025}, it needs no ground-truth depth supervision, no condition meta-text, and no condition annotations: competence is self-derived from the model's own calibrated per-modality predictions, using segmentation labels only. Finally, we report a systematic negative result on the remaining placement option, \emph{reliability as a pre-softmax attention-logit bias}. Nearby instantiations of that cell exist---a learned modality-prior additive bias for referring audio-visual segmentation~\cite{primed2026}, training-free entropy-derived additive attention modulation in LVLM decoding~\cite{sae2026}, and training-free multiplicative key scaling in SAM~2 memory attention for long videos~\cite{sam2long2024}---but none addresses multimodal dense segmentation. Our five-generation audit (Sec.~\ref{sec:experiments}) indicates why: once per-modality LoRA on a strong frozen backbone absorbs modality adaptation, logit-level suppression becomes a no-op, and reliability information is useful only at the gate and router.
